\section{Extending the household stage}\label{sec:roadmap}

Section~\ref{sec:results-secondstage} runs the energy and TCA food
episodes through PolicyEngine UK on the enhanced Family Resources
Survey. Extending the full statutory stage to the remaining three
episodes, and deepening the TCA run beyond its price vector, would add
episode-specific parameters to the same second stage; this study has
not implemented those runs, and four things would change when it is
done.

First, the consumption channel moves from decile tables to household
records: the enhanced FRS carries imputed household consumption across
the twelve COICOP divisions plus the sub-items that matter most for
trade shocks (petrol and diesel separately, in pounds and litres;
electricity and gas separately; bus fares), imputed by quantile
regression forest from the LCFS 2023--24 and calibrated to
administrative energy and fuel totals. Each episode's price vector
therefore prices each household's own basket, jointly with that
household's benefit entitlements, the covariance between exposure and
cushioning that decile tables cannot see. Tariff-line and
category-level first stages must be pre-aggregated to those roughly
sixteen categories, and the imputation freezes budget-share
\emph{patterns} at the 2023--24 donor vintage across simulated years, a
limitation the stage will state. Second, the earnings channels run through the companion study's
displacement machinery on the same households, so the two channels of a
given episode land on a common population and their incidence can be
compared household-by-household rather than distribution-by-distribution. Third, the rulebook axis would be scored directly: PolicyEngine's
parameter history carries dated values through the whole window, most
with a legislative or Commons Library reference (the Covid-era
Universal Credit uplift and its October 2021 removal, the March 2022
fuel-duty cut, the 2022 and 2023 upratings), its test suite exercises
the 2020--2023 rule years, and a counterfactual-rulebook run applies
an earlier year's parameters to a later survey vintage by copying the
survey inputs to that year, as Section~\ref{sec:results-rulebook}
does. Each episode
is therefore scored under the rules in force at its date, under the
preceding year's rules (isolating uprating lag), and with the
discretionary interventions switched off (isolating the automatic
component), the three-way decomposition of
Section~\ref{sec:framework} computed rather than approximated. Two
verified limits shape the earlier years: no 2021-vintage enhanced-FRS
data year exists (the earliest survey base is 2022--23), so 2021
rulebooks are applied to later households with incomes re-deflated;
and the survey's Universal Credit/legacy-benefit split follows
reported receipt rather than the historical rollout, so 2021--22
rule-year results understate legacy-benefit interactions and are
restricted to the tax and UC-parameter channels.

One methodological upgrade is specified here: the fixed-weight decile
indices of the first pass can be replaced by non-homothetic
distributional cost-of-living indices, which \citet{hochmuth2025}
construct from exactly the inputs this project already has (aggregate
prices, expenditure shares and a single cross-sectional expenditure
distribution), answering the standard objection that fixed baskets
misprice large relative-price movements.

Two limits will survive the full stage. The consumption
imputation carries its own error, inherited from the LCFS donor
models, and imputation error correlated with benefit status would bias
the exposure-cushioning covariance; the stage will report the
imputation diagnostics alongside the results. And the rulebook
decomposition attributes to ``discretion'' whatever the parameter
history records as legislated change, which bundles genuine shock
response with policy that would have happened anyway, a labelling
choice, stated as such, not an identification claim.
